% ============================================================
% 一页中文简历模板
% 编译：tectonic -X compile resume_template.tex
% 所有姓名 / 电话 / 邮箱 / 学校 / 项目链接均为占位符，请替换成自己的信息。
% 如不需要照片，当前文件可直接编译；需要照片时把 photo.jpg 放到同目录。
% photo.jpg 已加入 .gitignore，避免把真实照片提交到公共仓库。
% ============================================================
\documentclass[10pt, a4paper]{article}

\usepackage[top=0.8cm, bottom=0.4cm, left=1.3cm, right=1.3cm]{geometry}
% 使用 fandol 字体集，Tectonic / Overleaf / Linux 都能直接编译。
% 如果只在本机 Windows 上用 XeLaTeX，并且想用宋体/黑体，可改成 fontset=windows。
\usepackage[UTF8, fontset=fandol]{ctex}
\usepackage{xcolor}
\definecolor{primary}{RGB}{0, 82, 147}
\definecolor{secondary}{RGB}{100, 100, 100}
\usepackage{graphicx}
\usepackage{hyperref}
\hypersetup{
    colorlinks=true,
    linkcolor=primary,
    urlcolor=primary,
    pdfborder={0 0 0},
    pdftitle={Resume Template},
    pdfauthor={}
}
\usepackage{enumitem}
\setlist[itemize]{leftmargin=1.4em, nosep, topsep=0pt, itemsep=0pt}
\usepackage{fancyhdr}
\pagestyle{empty}
\usepackage{titlesec}
\titleformat{\section}{\normalsize\bfseries\color{primary}}{}{0em}{}[\titlerule]
\titlespacing{\section}{0pt}{5pt}{6pt}

% ========== 自定义命令 ==========
% 项目标题：项目名称（左）+ 时间（右）
\newcommand{\projectTitle}[2]{\noindent\textbf{#1} \hfill \textit{#2}\par\vspace{1pt}}
% 技术栈行
\newcommand{\projectTech}[1]{\noindent\textcolor{secondary}{\small 技术栈：#1}\par\vspace{1pt}}

\begin{document}

% ==================== 头部信息 ====================
\noindent
\begin{minipage}[c]{0.82\textwidth}
    {\LARGE \textbf{姓名}} \\[3pt]
    \textcolor{primary}{\textbf{求职意向：XXX 开发实习生}} \\[3pt]
    \small
    电话：1XX-XXXX-XXXX \quad | \quad 邮箱：\href{mailto:your@email.com}{your@email.com} \\
    \textbf{\textcolor{red}{可实习时间：X个月以上，可随时到岗（每周全勤5天）}}
\end{minipage}\hfill
\begin{minipage}[c]{0.15\textwidth}
    \raggedleft
    \IfFileExists{photo.jpg}{%
        \includegraphics[height=3.2cm]{photo.jpg}%
    }{%
        \rule{0pt}{3.2cm}%
    }
\end{minipage}

\vspace{-2pt}

% ==================== 教育背景 ====================
\section*{教育背景}
\noindent
\textbf{学校名称} \hfill 专业名称 \quad 硕士/本科 \hfill \textit{起始年月 - 结束年月}\par\vspace{2pt}
\noindent
\textbf{学校名称} \hfill 专业名称 \quad 硕士/本科 \hfill \textit{起始年月 - 结束年月}

% ==================== 项目经历 ====================
\section*{项目经历}

\projectTitle{\href{https://github.com/yourname/project1}{项目名称 — 一句话描述}}{202X.XX - 202X.XX}
\projectTech{技术栈1 + 技术栈2 + 技术栈3 + ...}
\noindent\textcolor{secondary}{\small \texttt{github.com/yourname/project1}}\par\vspace{1pt}
\begin{itemize}
    \item \textbf{小标题1}：具体做了什么，用了什么技术/方法，取得了什么可量化结果（\textbf{加粗关键数字}）。
    \item \textbf{小标题2}：具体做了什么，用了什么技术/方法，取得了什么可量化结果（\textbf{加粗关键数字}）。
    \item \textbf{小标题3}：具体做了什么，用了什么技术/方法，取得了什么可量化结果（\textbf{加粗关键数字}）。
    \item \textbf{小标题4}：具体做了什么，用了什么技术/方法，取得了什么可量化结果（\textbf{加粗关键数字}）。
    \item \textbf{小标题5}：具体做了什么，用了什么技术/方法，取得了什么可量化结果（\textbf{加粗关键数字}）。
\end{itemize}

\vspace{1pt}

\projectTitle{\href{https://github.com/yourname/project2}{项目名称 — 一句话描述}}{202X.XX - 至今}
\projectTech{技术栈1 + 技术栈2 + 技术栈3 + ...}
\noindent\textcolor{secondary}{\small \texttt{github.com/yourname/project2}}\par\vspace{1pt}
\begin{itemize}
    \item \textbf{小标题1}：具体做了什么，用了什么技术/方法，取得了什么可量化结果（\textbf{加粗关键数字}）。
    \item \textbf{小标题2}：具体做了什么，用了什么技术/方法，取得了什么可量化结果（\textbf{加粗关键数字}）。
    \item \textbf{小标题3}：具体做了什么，用了什么技术/方法，取得了什么可量化结果（\textbf{加粗关键数字}）。
    \item \textbf{小标题4}：具体做了什么，用了什么技术/方法，取得了什么可量化结果（\textbf{加粗关键数字}）。
\end{itemize}

% ==================== 专业技能 ====================
\section*{专业技能}
\begin{itemize}
    \item \textbf{计算机基础}：数据结构、操作系统、计算机网络基础扎实；熟悉 Linux 常用命令及 Git 工作流。
    \item \textbf{方向1}：熟悉 XXX 范式；熟悉 XXX 框架/协议；熟悉 XXX 机制。
    \item \textbf{方向2}：熟悉 XXX 异步 API；熟悉 XXX 数据库；熟悉 XXX 缓存/消息队列。
    \item \textbf{方向3}：熟悉 XXX 全链路开发；熟悉 XXX 工具/模型/策略。
    \item \textbf{工程化}：熟悉 Docker Compose 容器化部署；熟练使用 AI Coding 工具提升研发效率。
\end{itemize}

\end{document}
